% =====================================================================
\chapter{Funciones de varias variables, límites y continuidad}
\label{tema:3}
% =====================================================================

Una función de varias variables asigna un número a cada punto del plano o del espacio: la altitud de un terreno, la temperatura de una habitación o la pérdida de una red neuronal en función de sus parámetros. Este tema enseña a visualizarlas mediante gráficas y curvas de nivel y extiende a ellas las nociones de límite y continuidad. La novedad respecto a una variable es que un punto puede alcanzarse por infinitos caminos, y el límite debe ser el mismo por todos ellos. La presentación sigue las §4.1--4.2 de OpenStax~\cite{openstax3}, las sesiones~24--25 de 18.02SC~\cite{mit1802}, el capítulo~8 de Apostol~\cite{apostol2} y la §13.1 de Strang~\cite{strangcalc}.

\section{Funciones, dominio y gráfica}

\begin{definicion}[Función de varias variables]\label{def:funcion-n}
Una \emph{función real de $n$ variables} es una regla $f$ que asigna a cada punto $\vx=(x_1,\dots,x_n)$ de un conjunto $D\subseteq\R^n$, su \emph{dominio}, un único número real $f(\vx)$. Si no se indica, el dominio es el mayor conjunto en el que la fórmula tiene sentido. La \emph{imagen} es $\{f(\vx):\vx\in D\}$. Para $n=2$, la \emph{gráfica} de $f$ es la superficie $\{(x,y,z)\in\R^3: (x,y)\in D,\ z=f(x,y)\}$.
\end{definicion}

\addtocounter{definicion}{1}%

\section{Curvas y superficies de nivel}

\begin{definicion}[Conjuntos de nivel]\label{def:nivel}
Sea $f:D\subseteq\R^n\to\R$ y $c\in\R$. El \emph{conjunto de nivel} $c$ de $f$ es $\{\vx\in D: f(\vx)=c\}$. Para $n=2$ se llama \emph{curva de nivel}, y para $n=3$, \emph{superficie de nivel}. Un \emph{mapa de contorno} dibuja varias curvas de nivel, normalmente para valores de $c$ igualmente espaciados.
\end{definicion}

La curva de nivel $c$ es la proyección sobre el plano $xy$ del corte de la gráfica con el plano horizontal $z=c$, como las curvas de un mapa topográfico. Dos curvas de nivel distintas nunca se cortan, porque $f$ asigna un único valor a cada punto. Donde las curvas de nivel (equiespaciadas) están más juntas, la función cambia más deprisa.

\addtocounter{figure}{1}%

\addtocounter{definicion}{1}%

\section{Límites}

En $\R^n$ la cercanía se mide con la norma euclídea: la \emph{bola abierta} de centro $\va$ y radio $\delta>0$ es $B(\va,\delta)=\{\vx:\norm{\vx-\va}<\delta\}$. Se dice que $\va$ es un \emph{punto de acumulación} de $D$ si toda bola centrada en $\va$ contiene puntos de $D$ distintos de $\va$.

\begin{definicion}[Límite]\label{def:limite-n}
Sea $f:D\subseteq\R^n\to\R$ y $\va$ un punto de acumulación de $D$. Se dice que $\lim_{\vx\to\va}f(\vx)=L$ si para cada $\varepsilon>0$ existe $\delta>0$ tal que
\[
\vx\in D,\quad 0<\norm{\vx-\va}<\delta\ \Longrightarrow\ \abs{f(\vx)-L}<\varepsilon .
\]
\end{definicion}

La definición es la de una variable con $\abs{x-a}$ sustituido por $\norm{\vx-\va}$, y las demostraciones de las reglas algebraicas (límite de la suma, del producto, del cociente con denominador de límite no nulo) se trasladan sin cambios. Dos hechos elementales se usan constantemente: para $\vx=(x,y)$ y $\va=(a,b)$,
\begin{equation}\label{eq:t3-coordenada}
\abs{x-a}\le\norm{\vx-\va},\qquad \abs{y-b}\le\norm{\vx-\va},
\end{equation}
porque $\norm{\vx-\va}^2=(x-a)^2+(y-b)^2$ es mayor o igual que cada sumando.

\begin{proposicion}[Acotación]\label{prop:acotacion}
Si $\abs{f(\vx)-L}\le g(\vx)$ para $\vx$ cerca de $\va$ y $\lim_{\vx\to\va}g(\vx)=0$, entonces $\lim_{\vx\to\va}f(\vx)=L$.
\end{proposicion}

En el origen es útil pasar a \emph{coordenadas polares} $x=r\cos\theta$, $y=r\sin\theta$, con $r=\norm{(x,y)}$: si $\abs{f(r\cos\theta,r\sin\theta)-L}\le h(r)$ con $h(r)\to0$ cuando $r\to0^+$ y la cota $h$ \emph{no depende de $\theta$}, el límite es $L$ por la \cref{prop:acotacion}.

\addtocounter{definicion}{1}%

\begin{teorema}[Límites a lo largo de curvas]\label{teo:trayectorias}
Sea $\lim_{\vx\to\va}f(\vx)=L$ y sea $\boldsymbol\gamma:(0,\eta)\to D$ una curva con $\boldsymbol\gamma(t)\neq\va$ y $\lim_{t\to0^+}\boldsymbol\gamma(t)=\va$. Entonces $\lim_{t\to0^+}f(\boldsymbol\gamma(t))=L$.
\end{teorema}

\begin{corolario}[Criterio de no existencia]\label{cor:no-existe}
Si a lo largo de dos curvas que llegan a $\va$ la función tiene límites distintos (o a lo largo de una no tiene límite), entonces $\lim_{\vx\to\va}f(\vx)$ no existe.
\end{corolario}

\addtocounter{definicion}{1}%

\section{Continuidad}

\begin{definicion}[Continuidad]\label{def:continua-n}
Una función $f:D\to\R$ es \emph{continua en} $\va\in D$ si $\lim_{\vx\to\va}f(\vx)=f(\va)$ (o si $\va$ es un punto aislado de $D$). Es continua en $D$ si lo es en cada punto de $D$.
\end{definicion}

\begin{proposicion}[Funciones continuas elementales]\label{prop:continuas}
\begin{enumerate}
  \item Las proyecciones $(x,y)\mapsto x$ y $(x,y)\mapsto y$ son continuas en $\R^2$.
  \item Sumas, productos y cocientes (donde el denominador no se anula) de funciones continuas son continuos. En particular, los polinomios en $x$ e $y$ son continuos en $\R^2$, y las funciones racionales lo son en su dominio.
  \item Si $g$ es continua en $\va$ y $h:\R\to\R$ es continua en $g(\va)$, entonces $h\circ g$ es continua en $\va$.
\end{enumerate}
\end{proposicion}

Así, $e^{x^2+y^2}$, $\sin(xy)$ o $\ln(1+x^2+y^2)$ son continuas en todo $\R^2$, y para calcular su límite en un punto basta sustituir. Los casos interesantes son los cocientes en puntos donde se anula el denominador, como en los ejemplos anteriores.

\addtocounter{definicion}{1}%

\addtocounter{definicion}{1}%

\addtocounter{definicion}{1}%

